\section{What the fruit promises}

Augustine pictures a man who has confessed his sins and heard the word of God. He has good things in his possession and knows another person who needs them. One can tell what God has given this man when he takes pleasure in passing them on. Bellarmine pictures a field where no one has put a plough into the earth. Its fruit is the child Mary bears through the action of the Spirit. If we ask how the communicant can meet the Lord, Augustine helps us consider what a redeemed human being can learn to love. If we ask whom the communicant receives, Bellarmine helps us recognize the Son who took his humanity from Mary. A person can need the first change precisely because the second gift has been made to him.

There is a comparable difference between the fruit on the Gospel trees in Gregory and in Ambrose. For Gregory those who love God recognize the nearness of his kingdom as the old world's supports give way. For Ambrose the fruit on all the trees can be the confession of faith among many peoples, which will be followed by resurrection. His further warning about leaves that can cover sin prevents the hearer from treating every appearance of growth as evidence of faith. One Christian may await the end with particular longing because he has lately lost much that he loves. Another may have come to know Christ through the faith of people of another land. Both have reason to listen when Jesus tells the disciples whose coming they can expect.

After Communion the Church asks to prepare for the feast of her restoration. She has something to do in the time before it arrives, and someone whose birth she celebrates when it does. Paul's call to wake up addresses her conduct during that interval. His command to put on the Lord Jesus Christ names the One who can already be recognized in the conduct of those who keep his company.
